Language models told to keep a word secret still write stories from which a reader can tell which word they held. We ask whether this leak comes from planning and prose being entangled: a model that must make its own content decisions resolves them toward what it knows. If so, an explicit plan should remove the leak. We test this with \gemma as a white-box writer and \judge as a forced-choice (\twoafc) reader. A secret-free outline reduces discrimination from 90.0\% to 52.4\% (chance 50\%), and \gemmabig replicates the result (91.5\% to 51.5\%). A length-matched irrelevant passage leaves leakage unchanged (91.7\%), so the outline works by fixing content, not by diluting attention. An outline that the model writes while holding the secret does not help (87.6\%), because the leak moves into the plan: a secret-free writer following that outline still leaks (85.2\%). Plot twists behave the same way: openings leak a hidden twist (57.6\%), and a twist-free outline brings this to chance (49.3\%). Internally, the secret is decodable 15-way at 95--100\% at every story position in upper layers, even though the literal secret token is suppressed. An outline in context shrinks the secret's footprint at story tokens to about 30\%. Mean-ablating the secret subspace during generation cuts leakage to 60.5\%, against 80.7\% for a matched control subspace, and blocking attention to the secret sentence cuts it to 68.6\% with no loss in quality. To keep a secret out of generated text, the planner must not see it.
